\section{Conclusion}
\label{sec:conclusion}

Graph neural networks can effectively model habitat suitability for Swiss 
flora at 30\,m resolution. By encoding the landscape as a graph of
166,464 ecologically homogeneous patches linked by spatial
adjacency, the SOM-GNN-SDM framework captures spatial context
that traditional per-site models cannot represent.

Three main results:

\begin{enumerate}
  \item \textbf{Graph structure improves predictions.} The
    GNN-SDM outperforms the Random Forest baseline for
    85.3\,\% of 2,696 species (mean AUC 0.912 vs 0.874,
    mean TSS 0.583 vs 0.430), with the largest gains for
    species whose distributions depend on landscape
    connectivity rather than local features alone.

  \item \textbf{The model generalises to rare species.}
    Nationally threatened species achieve comparable or
    higher discrimination (AUC 0.874, TSS 0.708) than
    common species, confirming that the framework is
    suitable for conservation applications despite lower
    observation density.

  \item \textbf{Significant protection gaps exist.}
    Circuit-theory connectivity analysis across 241
    threatened species reveals that 41\,\% of
    high-priority habitat patches---locations that are both
    suitable habitat and critical connectivity
    corridors---currently lack any formal protection.
\end{enumerate}

These results should be interpreted as a methodological
proof-of-concept rather than a direct policy recommendation.
Conservation decisions involve many factors beyond ecology:
a landscape patch identified as a critical ecological corridor
may simultaneously be productive farmland, infrastructure, or
subject to competing land-use interests. The priority surface
identifies \emph{where} ecological connectivity is most
vulnerable, but the question of \emph{whether} and \emph{how}
to act requires integration with socio-economic constraints,
stakeholder consultation, and on-the-ground verification that
lie beyond the scope of this analysis.

The full pipeline---from raw satellite and climate data to
actionable conservation priority maps---runs end-to-end in
under 10 hours of compute time and is readily transferable to
other regions or taxa. As Switzerland continues to develop its
ecological infrastructure strategy, data-driven approaches like SOM-GNN-SDM 
can complement expert assessment and field surveys with spatially explicit, 
reproducible evidence.

\section*{Acknowledgements}

This work was conducted as part of the Certificate of Advanced
Studies in Applied Machine Learning at the University of Bern. 
Computational resources (AWS SageMaker GPU
instances) were provided by the Federal Office for the
Environment (FOEN).

This report does not constitute a policy recommendation of the
FOEN. While the author is employed at FOEN and the programme
was partially funded by FOEN, this work was carried out
independently as an academic case study exploring the
application of machine learning to conservation planning. Any
interpretation of results in a policy context would require
additional validation, stakeholder engagement, and integration
with existing planning instruments.

AI assistance was used in the creation of this report as well
as in the creation of the code used.

\paragraph{Code availability.}
The full analysis pipeline (notebooks 00--32) is available at
\url{https://github.com/kronmar-bafu/cas-aml-project}.
